\section{Plano de ação resumido}

O quadro reúne estratégia principal, ações-chave, recursos e indicadores, como pede a orientação [1]. A estratégia principal é o piloto territorial; a aquisição por confiança apoia a ativação.

\begin{table}[H]
\centering\small
\caption{Quadro-resumo do plano de ação.}
\begin{tabularx}{\textwidth}{P{2.4cm}YYY}
\toprule
\textbf{Estratégia} & \textbf{Ações-chave} & \textbf{Recursos} & \textbf{Indicadores} \\
\midrule
Piloto territorial (principal) & Entrevistas e mapeamento; revisão de créditos e QR; teste autônomo; piloto concierge de 8 semanas. & Horas da equipe; roteiros e termo de consentimento; ambiente de teste; contatos locais. & 8 de 10 tarefas concluídas sem ajuda; 80\% explicam créditos e QR; 2 clínicas aceitam novo ciclo. \\
Aquisição por confiança (apoio) & Conteúdo preventivo revisado; link por parceiro; funil até o primeiro lembrete. & Material simples; revisão por profissional qualificado; painel de métricas. & Ativação 60\%; retenção D30 30\%; custo por tutor ativado; NPS. \\
Escalabilidade & Padronizar ficha de parceiro; automatizar lembretes; definir regras de créditos. & Documentação de processos; back-end; revisão de segurança e LGPD. & Tempo de cadastro de parceiro; custo de suporte por tutor; margem por parceiro. \\
\bottomrule
\end{tabularx}
\end{table}

\begin{longtable}{P{1.8cm}P{4.6cm}P{4.2cm}P{2.6cm}}
\caption{Cronograma do ciclo de validação.}\label{tab:plano}\\
\toprule
\textbf{Período} & \textbf{Ação} & \textbf{Meta proposta} & \textbf{Responsável} \\
\midrule
\endfirsthead
\toprule
\textbf{Período} & \textbf{Ação} & \textbf{Meta proposta} & \textbf{Responsável} \\
\midrule
\endhead
Sem. 1--2 & Entrevistas e mapeamento local. & 15--20 tutores, 5 clínicas, 2 ONGs. & Giovanna e Kathleen \\
Sem. 3--4 & Revisão de créditos e QR; teste autônomo. & 10 testes, 8 concluídos sem ajuda; SUS aplicada. & Kathleen e Gabriel \\
Sem. 5--8 & Piloto concierge com registro de custos. & 40--50 tutores; 2--3 clínicas, se houver adesão. & Equipe; Júlia nos dados \\
Fechamento & Decidir entre iterar, continuar ou ampliar. & Ativação 60\%; D30 30\%; 2 clínicas. & Giovanna \\
\bottomrule
\end{longtable}

\subsection{Papéis na execução}

A divisão mantém os papéis da Atividade XI [30]: \textbf{Giovanna} coordena produto, entrevistas e parcerias; \textbf{Gabriel} cuida da arquitetura, estabilidade e segurança do aplicativo; \textbf{Júlia} responde por back-end, integrações e painel de métricas; \textbf{Kathleen} planeja os testes, sintetiza o feedback, revisa textos de interface e faz a ponte entre as frentes técnica e de produto.
